\chapter*{第 4 章\quad 流体动力学基本方程·孔珑题选}
\addcontentsline{toc}{chapter}{第 4 章\quad 流体动力学基本方程·孔珑题选}
\markboth{第 4 章\quad 流体动力学基本方程·孔珑题选}{}

\begin{tishi}
\textbf{题源：}本章题目取自孔珑主编《工程流体力学（第四版）》配套辅导书
《工程流体力学 学习指导及习题解答》（陈洁、袁铁江 编著，清华大学出版社）第 4 章
"流体运动学和流体动力学基础"的课后习题解答段（原书书页 42--58，原书题号 4-1 至 4-35）。
孔珑该章把\textbf{运动学与动力学合为一章}，故本册按西华 818 考纲（赵琴《流体力学与流体机械》第 1--8 章）
作了拆分：

\begin{tabularx}{\linewidth}{@{}lX@{}}
\toprule
原书题号 & 处理方式 \\
\midrule
4-13 $\sim$ 4-19、4-21 $\sim$ 4-31、4-34、4-35 & \textbf{本章收录}（流体动力学：伯努利方程、能量方程、动量方程、动量矩方程、泵与风机） \\
4-1 $\sim$ 4-12、4-20、4-32、4-33 & \textbf{本章不收}（只涉及速度场、流线/迹线、加速度、连续性方程、叶轮进出口流速与叶片速度三角形等运动学内容，归本丛书第 3 章"流体运动学"分册） \\
\bottomrule
\end{tabularx}

\textbf{做题要求：}① 伯努利方程类题必须写明\textbf{选断面（为什么选它）、基准面、压强基准（相对/绝对）、适用条件}
（不可压缩、恒定流、质量力只有重力、两点取在渐变流断面上）；
② 动量方程类题必须写明\textbf{控制体、外力分析、投影方向的正负号约定}，最后明确回答"是流体对壁面还是壁面对流体"；
③ 速度三角形类题必须写明\textbf{牵连速度、相对速度、绝对速度}三者的方向约定；
④ 全部题目都要有实质解答，不得写"略"。
原书题图一律用〔图示〕文字精确描述已知量，\textbf{不重绘图、不编造尺寸}。
凡 OCR 影印不清、题面数据只能由解答反推者，在章末按题号集中说明。
\end{tishi}

\begin{ti}\sloppy\emergencystretch=2em

\item[\textbf{4-13}] 如图 4-27 所示，直立圆管管径 $d_1=10\ \mathrm{mm}$，上端装有直径 $d_2=5\ \mathrm{mm}$ 的喷嘴，
喷嘴中心离圆管截面 1-1 的高度为 $h=3.6\ \mathrm{m}$，水从喷嘴向上排入大气，出口速度 $v_2=18\ \mathrm{m/s}$。
不计摩擦损失，计算截面 1-1 处所需要的计示压强 $p_{1e}$。
\tuyi{习题 4-27 图：一根铅垂圆管，下端为所取断面 1-1（管径 $d_1=10\ \mathrm{mm}$）；管上端接一段收缩喷嘴，喷嘴出口直径 $d_2=5\ \mathrm{mm}$，水自喷嘴向上射入大气（出口标有自由出流的短箭头），出口速度 $v_2=18\ \mathrm{m/s}$；喷嘴出口中心高出断面 1-1 的竖直距离为 $h=3.6\ \mathrm{m}$（以竖直尺寸线标出）。}
\xstarfull{4}\quad\yiju{E1 slide33「第四章 流体动力学基础＝考试计算题重点」；E2「第四章是主要考点、需掌握计算套路」；E6 教材 \S 4.3 总流伯努利方程（式 4.17、4.18）；本题为伯努利方程基本应用，真题无直接题号，故四星}\quad\nandu{中}

\item[\textbf{4-14}] 忽略损失，求图 4-28 所示文丘里管内的流量 $q_V$。
已知上游直管段直径 $d_1=300\ \mathrm{mm}$，喉部直径 $d_2=150\ \mathrm{mm}$，管内为水；
喉部与上游断面之间接一个倒置的\emph{空气}差压计（两臂上端充空气、下端与水管连通），两臂液面高差 $h=20\ \mathrm{cm}$。
\tuyi{习题 4-28 图：一段水平文丘里管，左端为上游直管段（断面 1-1，$d_1=300\ \mathrm{mm}$），中部收缩的喉部为断面 2-2（$d_2=150\ \mathrm{mm}$）；在断面 1-1 与 2-2 之间接一个倒置的 U 形差压计，两臂上端封入空气、下端与管内水连通，两臂液面高差 $h=20\ \mathrm{cm}$（图中以 20 cm 标注）；管中水流自左向右。}
\xstarfull{5}\quad\yiju{E4 2015 年计算题第 5 题「文丘里流量计」（同型）；E4 2026 回忆版计算第 2 题「文丘里管流量证明」；E1 slide33「第四章＝考试计算题重点」；E6 教材式 4.21 $q_V=K\sqrt{\Delta h}$}\quad\nandu{易}

\item[\textbf{4-15}] 图 4-29 所示为一文丘里管和压强计，试推导体积流量 $q_V$ 与压强计读数 $H$ 之间的关系式。
已知上游管径 $D$、喉部直径 $d$，管内为水（密度 $\rho$），压强计内为水银（密度 $\rho_m$），文丘里管水平放置。
\tuyi{习题 4-29 图：一段水平文丘里管（上游断面 1-1 直径 $D$、喉部断面 2-2 直径 $d$），在 1-1 与 2-2 之间接一支 U 形水银压强计，两臂液面高差为 $H$；图中另画出一条水平基准线，用以说明两断面的位置水头相同。}
\xstarfull{5}\quad\yiju{E4 2026 回忆版计算第 2 题「文丘里管流量证明」（推导型）；E4 2015 年计算题第 5 题「文丘里流量计」；E1 slide33「第四章＝考试计算题重点」；E6 教材式 4.19、4.20、4.21}\quad\nandu{难}

\item[\textbf{4-16}] 按图 4-30 所示的条件求当 $H=30\ \mathrm{cm}$ 时的流速 $u_1$。
图中压差计内液体的相对密度为 $0.8$，所测流体为水；两断面 1-1、2-2 等高，2-2 处为驻点（$u_2=0$）。
\tuyi{习题 4-30 图：一段输水管上引出两路测压管，与一个压差计相连，压差计内液体的相对密度为 0.8（比水轻）；所比较的两断面 1-1 与 2-2 在同一水平线上，其中断面 2-2 为驻点（$u_2=0$），压差计两臂液面高差 $H=30\ \mathrm{cm}$。}
\xstarfull{4}\quad\yiju{E4 2015 年填空题「U 形水银压差计」；E1 slide33「第四章＝考试计算题重点」；E6 教材式 4.11、4.12（压差计测点流速原理）；真题中该考点以填空形式出现，无计算题号，故四星}\quad\nandu{中}

\item[\textbf{4-17}] 输水管中水的计示压强为 $p_e=6.865\times10^5\ \mathrm{Pa}$。
假设法兰盘接头之间的填料破损，形成一个面积 $A=2\ \mathrm{mm^2}$ 的穿孔，
求该输水管一昼夜所漏损的水量。
\tuyi{习题 4-17 图：一段有压输水管，管中水的计示压强 $p_e=6.865\times10^5\ \mathrm{Pa}$（图中标注在管路上）；两个法兰盘接头之间的填料破损处形成一个小孔，孔口面积 $A=2\ \mathrm{mm^2}$，水自孔口向大气中射出。}
\xstarfull{4}\quad\yiju{E1 slide33「第四章＝考试计算题重点」、同行「第五章 孔口/管嘴/管路水力计算＝考试计算题重点」；E2「计算题占 50\% 分值」；E6 教材式 4.18（托里拆利关系 $u=\sqrt{2p/\rho}$）；本题实为小孔出流，真题无直接题号，故四星}\quad\nandu{易}

\item[\textbf{4-18}] 如图 4-31 所示，水池中的水沿一截面变化的水平管道排出，质量流量 $q_m=14\ \mathrm{kg/s}$。
若 $d_1=100\ \mathrm{mm}$，$d_2=75\ \mathrm{mm}$，$d_3=50\ \mathrm{mm}$，不计损失。
求所需的水头 $H$ 以及第二管段中央 $M$ 点的压强，并绘制测压管水头线。
\tuyi{习题 4-31 图：左侧为一个水池，池中自由液面通大气；水自池中沿一段水平管道向右排出，管道依次为 $d_1=100\ \mathrm{mm}$、$d_2=75\ \mathrm{mm}$、$d_3=50\ \mathrm{mm}$ 三段（第三段末端为自由出流，流入大气）；$M$ 点标在第二管段（$d_2=75\ \mathrm{mm}$）的中央；水池液面高出管道轴线 $H$（以竖直尺寸线标出）。}
\xstarfull{5}\quad\yiju{E4 2022 年 818 单选第 1 题「测压管水头线坡度小于零的沿程变化」；E4 2015 年填空题「变径管 $v_2$」；E1 slide33「第四章＝考试计算题重点」；E6 教材式 4.17、4.18 与"总水头线、测压管水头线"}\quad\nandu{中}

\item[\textbf{4-19}] 如图 4-32 所示，水从井 A 利用虹吸管引到井 B 中，设已知体积流量 $q_V=100\ \mathrm{m^3/h}$，
两井水面高差 $H=3\ \mathrm{m}$，虹吸管上端水平段高出井 B 出水管口 $z=6\ \mathrm{m}$，
不计虹吸管中的水头损失。试求虹吸管的管径 $d$ 及上端管中的负计示压强值 $p_e$。
\tuyi{习题 4-32 图：左侧为井 A、右侧为井 B，两井的水面高差 $H=3\ \mathrm{m}$（以竖直尺寸线标出）；一根虹吸管自井 A 伸入水下、向上翻越井壁，其\emph{上端水平段}位于两井水面之上，再向下弯折伸入井 B 的水中；虹吸管上端水平段高出井 B 出水管口 $z=6\ \mathrm{m}$（图中以 $z$ 标注）。}
\xstarfull{5}\quad\yiju{E4 2022 年 818 单选第 7 题「虹吸管最高处压强与大气压的关系」；E5 题库选择「虹吸管最高处压强小于大气压」；E1 slide33「第四章＝考试计算题重点」；E6 教材式 4.18 与虹吸管例题}\quad\nandu{中}

\item[\textbf{4-21}] 如图 4-33 所示，水沿渐缩管道垂直向上流动，已知 $d_1=30\ \mathrm{cm}$，$d_2=20\ \mathrm{cm}$，
计示压强 $p_{1e}=19.6\ \mathrm{N/cm^2}$，$p_{2e}=9.81\ \mathrm{N/cm^2}$，两断面高差 $h=2\ \mathrm{m}$。
若不计摩擦损失，试计算其流量 $q_V$。
\tuyi{习题 4-33 图：一段铅垂放置的渐缩管道，水自下向上流动；下断面 1-1 直径 $d_1=30\ \mathrm{cm}$、计示压强 $p_{1e}=19.6\ \mathrm{N/cm^2}$；上断面 2-2 直径 $d_2=20\ \mathrm{cm}$、计示压强 $p_{2e}=9.81\ \mathrm{N/cm^2}$；两断面之间的竖直高差为 $h=2\ \mathrm{m}$。}
\xstarfull{4}\quad\yiju{E1 slide33「第四章＝考试计算题重点」；E2「第四章是主要考点、需掌握计算套路」；E6 教材式 4.17、4.18 与连续性方程联立；真题无直接题号，故四星}\quad\nandu{中}

\item[\textbf{4-22}] 如图 4-34 所示，离心式水泵借一内径 $d=150\ \mathrm{mm}$ 的吸水管以 $q_V=60\ \mathrm{m^3/h}$ 的流量
从一口水槽中吸水，并将水送至压力水箱。设装在水泵与吸水管接头上的真空计指示出负压值 $39997\ \mathrm{Pa}$。
水力损失不计，试求水泵的吸水高度 $H_s$。
\tuyi{习题 4-34 图：左侧为一敞口盛水槽（液面通大气），一根内径 $d=150\ \mathrm{mm}$ 的吸水管自槽中伸入水下，向上接至离心式水泵；真空计装在吸水管与水泵的接头上，指示出负压值 39997 Pa；水泵轴线高出水槽液面 $H_s$（图中以 $H_s$ 标注）。}
\xstarfull{5}\quad\yiju{E5 题库选择「泵允许吸上真空高度 $H_s$」；E4 2022 年 818 单选第 7 题（负压与汽蚀判据同源）；E1 slide33「第四章＝考试计算题重点」；E6 教材式 4.17（水头损失与真空度）}\quad\nandu{中}

\item[\textbf{4-23}] 离心式风机借集流器从大气中吸取空气。其测压装置为一从直径 $d=200\ \mathrm{mm}$ 圆柱形管道上
接出的、下端插入水槽中的玻璃管。若水在玻璃管中的上升高度 $H=250\ \mathrm{mm}$，
求风机每秒钟吸取的空气量 $q_V$，空气的密度 $\rho=1.29\ \mathrm{kg/m^3}$（见图 4-35）。
\tuyi{习题 4-35 图：一段水平圆柱形管道（前端为风机的集流器进气口，管径 $d=200\ \mathrm{mm}$）；管壁上接出一根玻璃管，玻璃管下端插入一个水槽中；风机工作时管内形成负压，玻璃管中的水被吸上高度 $H=250\ \mathrm{mm}$；风机自左侧大气中吸取空气。}
\xstarfull{4}\quad\yiju{E1 slide33「第四章＝考试计算题重点」；E2「第四章需掌握计算套路」；E6 教材式 4.18 与毕托管（测压管）测速原理；真题无直接题号，故四星}\quad\nandu{中}

\item[\textbf{4-24}] 习题 4-24 图为一水平放置的 $90^\circ$ 收缩弯管，进口直径 $d_1=150\ \mathrm{mm}$、出口直径 $d_2=75\ \mathrm{mm}$；
弯管进口断面水的平均流速 $u_1=2.5\ \mathrm{m/s}$，进口计示静压强 $p_{1e}=6.86\times10^4\ \mathrm{Pa}$。
如不计能量损失，试求支撑弯管在其位置所需的水平力。
\tuyi{习题 4-24 图：一段水平放置的 $90^\circ$ 收缩弯管；进口断面 1-1 直径 $d_1=150\ \mathrm{mm}$，水流以平均流速 $u_1=2.5\ \mathrm{m/s}$ 沿管轴流入，进口计示静压强 $p_{1e}=6.86\times10^4\ \mathrm{Pa}$；出口断面 2-2 直径 $d_2=75\ \mathrm{mm}$，水流方向与进口方向成 $90^\circ$（均在水平面内）。}
\xstarfull{5}\quad\yiju{E4 2015 年计算题第 4 题「水平弯管动量方程求 $F_x$、$F_y$」（同型）；E1 slide33「第四章＝考试计算题重点」；E2「第四章是主要考点、需掌握计算套路」；E6 教材式 4.30、4.31}\quad\nandu{中}

\item[\textbf{4-25}] 额定流量 $q_m=35.69\ \mathrm{kg/s}$ 的过热蒸汽，压强 $p=981\ \mathrm{N/cm^2}$，温度 $t=510\ \mathrm{^\circ C}$
（对应的蒸汽比体积 $v=0.03067\ \mathrm{m^3/kg}$），经 $\phi 273\times23\ \mathrm{mm}$ 的主蒸汽管道铅垂向下，
再经 $90^\circ$ 弯管转向水平方向流动。如不计能量损失，试求蒸汽作用给弯管的水平力。
\tuyi{习题 4-25 图：一根 $\phi 273\times23\ \mathrm{mm}$ 的主蒸汽管道铅垂向下（外径 273 mm、壁厚 23 mm），蒸汽自上向下流动；管道末端接一个 $90^\circ$ 弯管，转向后蒸汽沿水平方向流动。}
\xstarfull{4}\quad\yiju{E1 slide33「第四章＝考试计算题重点」；E2「计算题占 50\% 分值」；E6 教材式 4.30、4.31（动量方程与压力项的处理）；真题未直接考气体/蒸汽动量，故四星}\quad\nandu{中}

\item[\textbf{4-26}] 如图 4-36 所示，相对密度为 $0.83$ 的油从直径 $d=50\ \mathrm{mm}$ 的喷嘴水平射向一直立平板，
射流速度 $u_0=20\ \mathrm{m/s}$，求支撑平板所需的力 $F$。
\tuyi{习题 4-36 图：一根水平的油射流管自左侧向右射出一股油射流（相对密度 0.83、喷嘴直径 $d=50\ \mathrm{mm}$、速度 $u_0=20\ \mathrm{m/s}$）；射流垂直冲击一块直立平板（图中以竖直双线示意），平板由支撑力 $F$ 固定。}
\xstarfull{4}\quad\yiju{E1 slide33「第四章＝考试计算题重点」；E2「计算题占 50\% 分值、需掌握计算套路」；E4 2026 回忆版计算第 3 题为「书上动量方程例题」；E6 教材式 4.30、4.31 与例 4.4（射流冲击平板 $F=\rho q_V u$）；真题未把射流冲击平板单独成题，故四星}\quad\nandu{易}

\item[\textbf{4-27}] 标准状态的空气从喷嘴里射出，吹到一与之成直角的壁面上，
壁面上装有测压计，测压计读数高于大气压 $466.6\ \mathrm{Pa}$。求空气离开喷嘴时的速度。
\tuyi{习题 4-27 图：一个喷嘴向右侧射出一股空气射流；射流垂直（成直角）冲击一块壁面，壁面上开有测压孔并接测压计，测压计读数高于大气压 466.6 Pa。}
\xstarfull{3}\quad\yiju{E6 教材式 4.30、4.31 与滞止压强（动压）概念；E1 slide33「第四章＝考试计算题重点」；真题未出现该题，属基本公式套用，故三星}\quad\nandu{中}

\item[\textbf{4-28}] 水龙头与压力箱连接，计示压强 $170\ \mathrm{kPa}$。水龙头入口处的流速可忽略不计，
水龙头往大气中喷水，设水柱成单根流线、忽略空气阻力，估计水流离出口能达到的最大高度。
\tuyi{习题 4-28 图：一个压力箱与一个竖直向上的水龙头相连，箱内水的计示压强为 170 kPa；水龙头向大气中喷出一股铅垂水柱，水柱最高点高出水龙头出口的高度记为 $H$。}
\xstarfull{3}\quad\yiju{E6 教材式 4.18（压强水头与高度水头互化）；E1 slide33「第四章＝考试计算题重点」；真题未直接出现，属基本公式套用，故三星}\quad\nandu{易}

\item[\textbf{4-29}] 如图 4-37 所示，一股射流以速度 $u_0$ 水平射到倾斜光滑平板上，体积流量为 $q_{V0}$，
平板与射流轴线的夹角为 $\theta$。求沿板面向两侧的分流流量 $q_{V1}$ 与 $q_{V2}$ 的表达式，
以及流体对板面的作用力。忽略流体撞击的损失和重力影响，射流的压强分布在分流前后都没有变化。
\tuyi{习题 4-37 图：一股速度为 $u_0$、体积流量为 $q_{V0}$ 的水平射流自左侧射到一块倾斜光滑平板上，平板与射流轴线的夹角为 $\theta$；射流在板上分成两股沿板面流走：沿板面向上的一股流量为 $q_{V1}$，沿板面向下的一股流量为 $q_{V2}$；两股分流与来流的压强均为大气压。}
\xstarfull{4}\quad\yiju{E1 slide33「第四章＝考试计算题重点」；E2「第四章是主要考点、需掌握计算套路」；E6 教材式 4.30、4.31 与例 4.4（射流冲击斜置平板）；真题未直接成题，故四星}\quad\nandu{难}

\item[\textbf{4-30}] 如图 4-37 所示的流动，如果沿下侧流动的流量为流体总流量的 45\%，
问平板倾斜角 $\theta$ 多大？
\tuyi{同习题 4-37 图：水平射流以 $u_0$ 射到与射流轴线成 $\theta$ 角的倾斜光滑平板上，分成沿板面向上的 $q_{V1}$ 与沿板面向下的 $q_{V2}$ 两股；本题给定 $q_{V2}/q_{V0}=45\%$。}
\xstarfull{3}\quad\yiju{E6 与第 4-29 题同一结论（式 4.30、4.31 的推论）；E1 slide33「第四章＝考试计算题重点」；真题未出现，属既有公式的反用，故三星}\quad\nandu{易}

\item[\textbf{4-31}] 如图 4-38 所示，平板向着射流以等速 $v$ 运动，射流速度为 $u_0$、喷嘴面积为 $A_0$、
平板与射流轴线夹角为 $\theta$，导出使平板运动所需功率的表达式。
\tuyi{习题 4-38 图：一股速度为 $u_0$ 的水平射流自喷嘴（面积 $A_0$）射出；一块倾斜平板（与射流轴线成 $\theta$ 角）以等速 $v$ 迎着射流运动（图中以指向射流方向的箭头标注平板运动方向）。}
\xstarfull{2}\quad\yiju{E6 教材式 4.30、4.31 的导出应用（相对运动与动量方程结合）；E1 slide33「第四章＝考试计算题重点」；真题未出现且超出基本套路，属扩展题，故二星}\quad\nandu{难}

\item[\textbf{4-34}] 图 4-41 中的风机叶轮的内径 $d_1=12.5\ \mathrm{cm}$，外径 $d_2=30\ \mathrm{cm}$，叶片宽度 $b=2.5\ \mathrm{cm}$，
转速 $n=1725\ \mathrm{r/min}$，体积流量 $q_V=372\ \mathrm{m^3/h}$。空气在叶片进口处沿径向流入，
绝对压强 $p_1=9.7\times10^4\ \mathrm{Pa}$，气温 $t=20\ \mathrm{^\circ C}$，叶片出口方向与叶轮外缘切线方向的夹角 $\beta_2=30^\circ$。
假设流体是理想不可压缩流体：
\begin{xiaowen}
\item 画出进口处的速度图，并计算叶片的进口角 $\beta_1$；
\item 画出出口处的速度图，并计算出口处的绝对速度 $v_2$；
\item 求所需的扭矩 $M$。
\end{xiaowen}
\tuyi{习题 4-41 图：一个离心式风机叶轮（内径 $d_1=12.5\ \mathrm{cm}$、外径 $d_2=30\ \mathrm{cm}$、叶片宽度 $b=2.5\ \mathrm{cm}$，转速 $n=1725\ \mathrm{r/min}$）；图中给出叶轮进、出口的速度三角形：进口绝对速度沿径向（$\alpha_1=90^\circ$），出口叶片方向与叶轮外缘切线的夹角 $\beta_2=30^\circ$。}
\xstarfull{4}\quad\yiju{E6 教材式 4.32、4.33 与例 4.5（离心风机叶轮：题面数据与本例相同，$n=1725\ \mathrm{r/min}$、$d_1=125\ \mathrm{mm}$、$d_2=300\ \mathrm{mm}$、$b=25\ \mathrm{mm}$、$q_V=372\ \mathrm{m^3/h}$、$\beta_2=30^\circ$）；E1 slide33「第四章＝考试计算题重点」；动量矩方程真题未独立成题，故四星}\quad\nandu{难}

\item[\textbf{4-35}] 图 4-42 所示为一有对称臂的洒水器，设总体积流量为 $5.6\times10^{-4}\ \mathrm{m^3/s}$，
每个喷嘴的出口面积为 $0.93\ \mathrm{cm^2}$（两喷嘴共 $2\times0.93\ \mathrm{cm^2}$），转臂长 $r=0.25\ \mathrm{m}$，喷嘴轴线与转臂成 $45^\circ$。
不计摩擦，求它的角速度；如不让它转动，应施加多大扭矩。
\tuyi{习题 4-42 图：一个有对称臂的洒水器，两臂在同一水平面内、绕中心的铅垂轴转动；每个转臂长 $r=0.25\ \mathrm{m}$，臂端喷嘴的喷水方向与转臂成 $45^\circ$（图中以 $45^\circ$ 角标注）；水从两喷嘴喷出，图中以弧形箭头表示转动方向与角速度 $\omega$。}
\xstarfull{4}\quad\yiju{E6 教材式 4.32（动量矩方程）与 \S 4.4 叶轮机械基本方程；E1 slide33「第四章＝考试计算题重点」；E2「第四章是主要考点、需掌握计算套路」；动量矩方程真题未独立成题，故四星}\quad\nandu{难}

\end{ti}

\begin{tishi}
\textbf{〔收录与不收录的对照〕}孔珑该章 35 题中，本章收录 20 题（4-13、4-14、4-15、4-16、4-17、4-18、4-19、4-21、4-22、4-23、
4-24、4-25、4-26、4-27、4-28、4-29、4-30、4-31、4-34、4-35），
\textbf{不收 15 题}：4-1（绕圆柱平面流动求驻点与速度分布）、4-2（求流线方程）、4-3（两平行平板间速度分布）、
4-4、4-5、4-6（判断维数与求加速度）、4-7、4-8、4-9、4-10、4-11、4-12（连续性方程求流速、流量、质量流量与管径）、
4-20（送风管道质量流量、体积流量与平均流速，只用到连续性方程与状态方程）、
4-32（射流流入叶片，仅由速度三角形求叶片出口边角度）、4-33（水泵叶轮进出口流速，仅用到连续性方程）。
上述 15 题均属\textbf{流体运动学}，归本丛书第 3 章"流体运动学"分册。

\textbf{〔影印不清处（解题时请以题图为准）〕}
\textbf{第 4-14 题：}原书图 4-28 的图注影印不清，上游直管段直径 $d_1=300\ \mathrm{mm}$ 由原书解答中
$u_2=4u_1$ 与 $q_V=0.0361\ \mathrm{m^3/s}$ 反推得到（喉部直径 15 cm、读数 20 cm 为图注原有）；
若你的题图另有上游直径，请按图取值，方法与解答完全相同。\\
\textbf{第 4-16 题：}原书图 4-30 的管路与压差计接法影印不清，本章按原书解答反推出的条件表述为
"两断面等高、2-2 为驻点、压差计内液体相对密度 0.8、读数 $H=30\ \mathrm{cm}$"；
若你的题图接法不同，请按"由压差计读数求流速"的同一方法处理。\\
\textbf{第 4-24 题：}原书该题题面开头的管径标注在 OCR 中脱漏，进口直径 $d_1=150\ \mathrm{mm}$ 由原书解答
（连续性方程、$x$ 向动量方程中的 $\frac{\pi d_1^2}{4}$ 数值、出口流速 $u_2=10\ \mathrm{m/s}$）反推得到；
$d_2=75\ \mathrm{mm}$、$u_1=2.5\ \mathrm{m/s}$、$p_{1e}=6.86\times10^4\ \mathrm{Pa}$ 为原书原有数据。\\
\textbf{第 4-26 题：}喷嘴直径 $d=50\ \mathrm{mm}$ 仅在原书解答中出现（$\frac{\pi(0.05)^2}{4}$），题面未重复给出，请按题图核对。\\
\textbf{第 4-27 题：}原书按动量方程给出 $p=\rho u^2$（壁面\emph{平均}压强）的算法，答案 $u\approx19\ \mathrm{m/s}$；
若把测压计读数理解为驻点\emph{总压} $\frac{1}{2}\rho u^2$，则 $u\approx26.9\ \mathrm{m/s}$。两种口径的差别与取舍见解答后的说明。\\
\textbf{第 4-35 题：}转臂长 $r=0.25\ \mathrm{m}$ 与喷嘴倾角 $45^\circ$ 由原书解答中的
$v_{2e}=\omega r=0.25\omega$、$v_2\sin45^\circ$ 反推得到；若题图另有尺寸，请按图取值，方法不变。
\end{tishi}
